\section{ADDITIONAL END-TO-END EXAMPLES}
\label{app:examples}

The figures below complement Figure~\ref{fig:qualitative} with two further
source/target page pairs produced by the complete
\texttt{parse}\,$\to$\,\texttt{translate}\,$\to$\,\texttt{render} pipeline. In
every case, images and \texttt{BYPASS} regions are composited unchanged from the
source page, and the translated text is typeset inside the original bounding
boxes at a font size chosen by the group-consistent sizing rule
(Section~\ref{subsubsec:sizing}). Figure~\ref{fig:app-slide-2} shows a slide
and Figure~\ref{fig:app-book-2} a book page.

\begin{figure}[H]
  \centering
  \begin{subfigure}[t]{0.46\textwidth}
    \centering
    \includegraphics[width=\linewidth]{source_2.png}
    \caption{Source page}
  \end{subfigure}\hfill
  \begin{subfigure}[t]{0.46\textwidth}
    \centering
    \includegraphics[width=\linewidth]{des_2.png}
    \caption{Translated page}
  \end{subfigure}
  \caption{End-to-end result on a presentation-slide PDF. Short boxed text is
  the hardest case for the length constraint, since each box has little slack;
  the group-wise sizing rule (Section~\ref{subsubsec:sizing}) keeps sibling boxes
  visually consistent, and the per-element appearance estimate
  (Section~\ref{subsubsec:colour}) keeps the coloured fills and non-black type of the
  source slide instead of masking them white.}
  \label{fig:app-slide-2}
\end{figure}

\begin{figure}[H]
  \centering
  \begin{subfigure}[t]{0.46\textwidth}
    \centering
    \includegraphics[width=\linewidth]{source_4.png}
    \caption{Source page}
  \end{subfigure}\hfill
  \begin{subfigure}[t]{0.46\textwidth}
    \centering
    \includegraphics[width=\linewidth]{des_4.png}
    \caption{Translated page}
  \end{subfigure}
  \caption{End-to-end result on a book page with dense flowing text. Because
  paragraphs have vertical slack, overflow rarely collides with a neighbour and
  the canonical group font size is retained.}
  \label{fig:app-book-2}
\end{figure}
